% ===========================================================================
%  第15章：级联姿态-角速度控制架构
% ===========================================================================

\section{级联姿态-角速度控制架构：解耦设计与交叉项补偿}
\label{sec:cascade_chapter}

前文分别发展了三种姿态控制架构：第\ref{sec:control}节的级联SAS（对角分配）、
第\ref{sec:INDI}节的INDI（在线效能矩阵+角加速度反馈）、
第\ref{sec:lqr}节的LQR（工作点最优全状态反馈）。
第\ref{sec:coupling_rejection}节与第\ref{sec:lqr_coupling_limit}节的耦合抑制对比表明，
三者在处理本构型固有的反扭矩投影交叉耦合时各有局限：
SAS的对角分配结构性放弃交叉项，INDI受角加速度估计噪声与增量饱和制约，
LQR的冻结效能矩阵随工作点漂移而退化。
本章提出一种级联姿态-角速度控制架构，通过时间尺度分离将姿态控制与交叉项补偿解耦，
在保留各方案优势的同时规避其固有缺陷，并给出完整的理论分析与仿真验证。

\subsection{设计思想：时间尺度分离与功能解耦}
\label{sec:cascade_philosophy}

单级控制架构的根本困难在于：姿态误差动力学、角速度动力学与执行器交叉耦合
被压缩在同一个反馈回路中，控制律必须同时兼顾跟踪品质、稳定性与耦合补偿，
三者相互掣肘。第\ref{sec:lqr_coupling_limit}节的仿真表明，即使将LQR的效能矩阵
在线化（增益调度、SDRE或在线分配），只要保留单级全状态反馈结构，
交叉项补偿就难以突破$1^\circ$量级——补偿力矩被姿态-角速度耦合动态"淹没"。

级联架构的设计思想是\textbf{用时间尺度分离换取功能解耦}：
将控制问题分解为两个带宽分离的回路，每个回路只承担单一职能。

\begin{itemize}
    \item \textbf{外环（姿态环，慢环）}：仅负责姿态误差的调节，输出角速度指令。
          由于不直接操作执行器，外环无需考虑效能矩阵、执行器饱和与交叉耦合，
          可采用极简的纯比例（P）或比例-积分（PI）结构；
    \item \textbf{内环（角速度环，快环）}：仅负责跟踪角速度指令，输出执行器指令。
          由于角速度环带宽显著高于姿态环（时间尺度分离假设），
          内环可在姿态还未来得及响应前快速抑制交叉耦合扰动，
          使其无法经角速度-姿态积分传播到姿态环。
\end{itemize}

这一思想的物理本质是：\textbf{交叉耦合的抑制速度必须快于它影响姿态的速度}。
在单级架构中，补偿力矩与耦合扰动同时作用于同一动力学，无法形成时间尺度上的先后；
级联架构则通过内环的高带宽，使耦合扰动在姿态环"察觉"之前即被消除。
这与第\ref{sec:control}节SAS的级联思想一脉相承，但内环用在线效能矩阵
替代了对角近似，从而把几何解耦升级为精确解耦。

\subsection{控制律设计}
\label{sec:cascade_design}

\subsubsection{外环：误差四元数比例-积分控制}

姿态反馈采用第\ref{sec:quat_control}节的误差四元数，
以规避欧拉角在$\theta=\pm90^\circ$处的万向锁奇异性
（对第\ref{sec:vtol}节的tail-sitter悬停至关重要）。
设当前姿态四元数$\mathbf{q}$、期望姿态$\mathbf{q}_{\mathrm{des}}$，
误差四元数$\mathbf{q}_{\mathrm{err}}=\mathbf{q}_{\mathrm{des}}^{*}\otimes\mathbf{q}$
的矢量部分$\boldsymbol{\varepsilon}_{\mathrm{err}}$编码了姿态误差的方向与大小。
外环将其映射为角速度指令：

\begin{equation}
    \boxed{
    \boldsymbol{\omega}_{\mathrm{ref}}
    = -\mathbf{K}_p\,\boldsymbol{\varepsilon}_{\mathrm{err}}
      - \mathbf{K}_i\int\boldsymbol{\varepsilon}_{\mathrm{err}}\,\mathrm{d}t
      + \boldsymbol{\omega}_{\mathrm{ref}}^{\mathrm{ff}}
    }
    \label{eq:cascade_outer_ch}
\end{equation}

其中$\mathbf{K}_p$、$\mathbf{K}_i$为正定对角增益矩阵，
$\boldsymbol{\omega}_{\mathrm{ref}}^{\mathrm{ff}}$为期望轨迹的角速度前馈
（如slerp过渡段，见第\ref{sec:quat_control}节）。
在小姿态误差下$\boldsymbol{\varepsilon}_{\mathrm{err}}\approx\frac{1}{2}\Delta$姿态角，
故$\mathbf{K}_p$的物理意义是姿态回正的特征速率。
积分项用于消除常值扰动（如残余推力矩、风载）引起的稳态误差；
纯比例外环在常值耦合扰动下存在稳态偏差（仿真中约$0.04^\circ$），
积分项可将其进一步压缩。

\subsubsection{内环：角速度环与在线效能矩阵分配}

内环将角速度误差映射为虚拟角加速度，再经在线效能矩阵分配为执行器增量：

\begin{align}
    \boldsymbol{\nu} &= \mathbf{K}_\omega\left(\boldsymbol{\omega}_{\mathrm{ref}}
                        - \boldsymbol{\omega}\right)
    \label{eq:cascade_nu_ch}\\
    \mathbf{u}_{k+1} &= \mathbf{u}_k
        + \mathbf{B}_{\mathrm{true}}^{-1}(\mathbf{u}_k)\,\boldsymbol{\nu}
    \label{eq:cascade_alloc_ch}
\end{align}

其中$\mathbf{u}=[\Delta\omega,\delta_t,\delta_f]^T$，
$\mathbf{B}_{\mathrm{true}}(\mathbf{u}_k)$为式(\ref{eq:B_true})的在线Jacobian，
$\mathbf{K}_\omega$为角速度环比例增益。
内环的两个关键设计：

\begin{enumerate}[label=(\arabic*)]
    \item \textbf{交叉项精确补偿}：$\mathbf{B}_{\mathrm{true}}^{-1}$的非对角元
          将前摆对俯仰的反扭矩耦合$-\tau_f\cos\delta_f$、
          尾摆对偏航的耦合$-\tau_t\cos\delta_t$显式计入分配，
          实现任意工作点下的精确解耦（而非LQR的工作点冻结或SAS的对角近似）；
    \item \textbf{无角加速度估计}：虚拟角加速度$\boldsymbol{\nu}$
          由角速度误差经比例生成（式\ref{eq:cascade_nu_ch}），
          $\mathbf{B}_{\mathrm{true}}^{-1}$仅作静态分配而非动态逆。
          这与INDI的本质区别在于：INDI需要实测$\dot{\boldsymbol{\omega}}$
          来构造增量逆，而级联内环不需要——从而彻底规避了
          第\ref{sec:estimator_study}节分析的角加速度估计噪声放大问题。
\end{enumerate}

\subsection{理论分析}
\label{sec:cascade_theory}

\subsubsection{与INDI的本质区别}

级联内环（式\ref{eq:cascade_alloc_ch}）与INDI（式\ref{eq:indi_our_config}）
形式上都含$\mathbf{B}_{\mathrm{true}}^{-1}$，但控制律结构根本不同：

\begin{table}[H]
\centering
\caption{级联内环与INDI的结构对比}
\label{tab:cascade_vs_indi}
\small
\begin{tabularx}{\textwidth}{l X X}
\toprule
 & \textbf{级联内环} & \textbf{INDI} \\
\midrule
反馈量 & 角速度$\boldsymbol{\omega}$ & 角加速度$\dot{\boldsymbol{\omega}}$（实测/估计） \\
控制律 & $\mathbf{u}_{k+1}=\mathbf{u}_k+\mathbf{B}^{-1}\mathbf{K}_\omega(\boldsymbol{\omega}_{\mathrm{ref}}-\boldsymbol{\omega})$ & $\mathbf{u}_{k+1}=\mathbf{u}_k+\mathbf{B}^{-1}(\boldsymbol{\nu}-\dot{\boldsymbol{\omega}}_{\mathrm{IMU}})$ \\
动态逆 & 否（仅静态分配） & 是（增量动态逆） \\
噪声敏感 & 低（$\boldsymbol{\omega}$直接可测） & 高（$\dot{\boldsymbol{\omega}}$需差分/滤波） \\
未建模动态对消 & 依赖反馈带宽 & 依赖$\dot{\boldsymbol{\omega}}$隐式对消 \\
\bottomrule
\end{tabularx}
\end{table}

级联内环放弃了INDI"用实测角加速度对消未建模动态"的能力，
换取了对角加速度估计噪声的免疫。这一取舍在本构型上是合理的：
第\ref{sec:propulsion}节已分析陀螺效应在前后转子间一阶抵消，
Coriolis耦合在小角速率下为高阶小量，故未建模动态本就较弱，
INDI的对消能力优势有限，而其噪声敏感性却是实质缺陷。

\subsubsection{稳定性与时间尺度分离条件}

级联架构的稳定性依赖内外环的时间尺度分离。设外环带宽$\omega_{\mathrm{BW,o}}$、
内环带宽$\omega_{\mathrm{BW,i}}$，则要求
\begin{equation}
    \omega_{\mathrm{BW,i}} \gtrsim 3\text{--}5\,\omega_{\mathrm{BW,o}}
    \label{eq:tss_condition}
\end{equation}
以保证内环在姿态环响应前完成角速度跟踪。外环纯比例的带宽约为
$\omega_{\mathrm{BW,o}}\approx K_p/2$（因$\dot{\boldsymbol{\varepsilon}}\approx\frac{1}{2}\boldsymbol{\omega}$），
内环带宽约为$\omega_{\mathrm{BW,i}}\approx K_\omega$。
仿真采用$K_p=2.0$、$K_\omega=8.0$，得
$\omega_{\mathrm{BW,o}}\approx1$\,rad/s、$\omega_{\mathrm{BW,i}}\approx8$\,rad/s，
分离比约8倍，满足式(\ref{eq:tss_condition})且留有裕度。

需要强调：时间尺度分离是\textbf{近似}，严格稳定性需经Lyapunov分析或
奇异摄动理论验证。在概念阶段，本架构的稳定性由仿真时域响应间接确认；
工程实现前应补充严格的稳定性证明与裕度评估。

\subsubsection{计算复杂度}

内环每周期需一次$3\times3$矩阵在线重构与求逆
（$\mathbf{B}_{\mathrm{true}}$由当前$\sin/\cos$摆角与推力/反扭矩组合而成，
伴随矩阵法求逆约18次乘法+1次除法），计算量与INDI相当，
远低于SDRE的每周期CARE求解（$6\times6$ Riccati方程）。
外环仅向量比例-积分运算。整体计算开销在嵌入式飞控平台上可忽略。

\subsection{仿真验证}
\label{sec:cascade_sim}

\subsubsection{巡航工况交叉耦合抑制}

在24\,m/s巡航配平、前摆$\delta_f$阶跃产生俯仰耦合的场景下
（设置同第\ref{sec:coupling_rejection}节），级联架构与LQR、INDI的
俯仰耦合峰值对比如表\ref{tab:cascade_ch_coupling}所示。

\begin{table}[H]
\centering
\caption{级联架构与LQR/INDI的俯仰耦合峰值对比（°，5次蒙特卡洛均值）}
\label{tab:cascade_ch_coupling}
\small
\begin{tabular}{l c c}
\toprule
\textbf{控制架构} & $\delta_f=10^\circ$ & $\delta_f=25^\circ$ \\
\midrule
LQR（冻结$\mathbf{B}$，单级） & 1.33 & 2.02 \\
INDI（在线$\mathbf{B}$，单级） & 0.40 & 2.86 \\
\textbf{级联P外环 + 在线$\mathbf{B}_{\mathrm{true}}$内环} & \textbf{0.02} & \textbf{0.04} \\
\bottomrule
\end{tabular}
\end{table}

级联架构将俯仰耦合峰值从LQR的$1.33$--$2.02^\circ$、INDI的$0.40$--$2.86^\circ$
降至$0.02$--$0.04^\circ$，改善约两个数量级。
时域上，$\delta_f=25^\circ$阶跃后俯仰角仅有$0.047^\circ$的微小瞬态即恢复，
而LQR持续偏离$2.1^\circ$。纯比例外环存在约$0.04^\circ$稳态误差，
引入积分项后进一步压缩。

\subsubsection{Tail-sitter大姿态角验证（万向锁附近）}

为验证误差四元数外环在万向锁附近的有效性，在tail-sitter垂直起飞工况
（3秒slerp过渡，水平$\to$悬停$\theta=-90^\circ$）下对比四种架构。
结果见表\ref{tab:cascade_ch_vtol}。

\begin{table}[H]
\centering
\caption{Tail-sitter垂直起飞姿态误差对比（3\,s slerp）}
\label{tab:cascade_ch_vtol}
\small
\begin{tabular}{l c c}
\toprule
\textbf{控制架构} & $\|\boldsymbol{\varepsilon}_{\mathrm{err}}\|$峰值 & 末值（10\,s） \\
\midrule
SAS（增益调度+前馈） & 0.198 & $4.4\times10^{-2}$ \\
INDI（v2） & 0.067 & $1.5\times10^{-3}$ \\
LQR & 0.017 & $2.1\times10^{-5}$ \\
\textbf{级联（误差四元数外环）} & 0.028 & $3.9\times10^{-3}$ \\
\bottomrule
\end{tabular}
\end{table}

级联架构在万向锁附近（$\theta=-90^\circ$）工作正常，
姿态误差峰值$0.028$（对应约$3^\circ$），远优于SAS的$0.198$，
接近LQR的$0.017$。误差四元数外环在欧拉角奇异的悬停姿态下保持正则，
验证了第\ref{sec:quat_control}节四元数反馈对级联外环的必要性。
其末值误差（$3.9\times10^{-3}$）介于INDI与LQR之间，
可通过增大积分增益或前馈进一步优化。

\subsection{分析讨论}
\label{sec:cascade_analysis}

\subsubsection{三种"LQR在线B"路径的失败教训}

在探索LQR交叉项补偿在线化时，曾系统评估三种直接路径
（增益调度LQR、SDRE在线Riccati、LQR外环+在线B分配，
实现见\texttt{lqr\_online\_b.py}）。仿真表明三者的耦合抑制均未能显著改善
（峰值$1.3$--$2.6^\circ$，与纯LQR相当）。根本原因是它们都保留了
\textbf{单级全状态反馈}结构——无论效能矩阵如何更新，
交叉项补偿仍与姿态-角速度耦合动态竞争。
这一失败从反面印证了级联架构的核心论点：
\textbf{交叉项补偿的有效性更多取决于控制架构的解耦程度，
而非效能矩阵的在线化程度}。

\subsubsection{架构的适用边界}

级联架构并非万能，其有效性受以下条件约束：

\begin{enumerate}[label=(\arabic*)]
    \item \textbf{时间尺度分离}：内环带宽必须足够高。当执行器响应迟缓
          （大惯量负载、低带宽舵机）使内环带宽受限时，解耦假设退化，
          耦合抑制能力下降；
    \item \textbf{低油门退化}：与INDI相同，$\omega_0\to0$时
          $\mathbf{B}_{\mathrm{true}}$趋于奇异，内环分配失效。
          需在第\ref{sec:INDI}节所述的低油门阈值处回退到对角分配或舵面控制；
    \item \textbf{执行器饱和}：在线$\mathbf{B}_{\mathrm{true}}$的补偿指令
          受摆角/差速限幅约束，大扰动下饱和会使补偿失能
          （这是INDI在$\delta_f=25^\circ$反超的原因，级联通过完整力矩分配
          而非增量限幅缓解，但未根本消除）。
\end{enumerate}

\subsubsection{与其他架构的定位关系}

级联架构并非取代LQR或INDI，而是在特定需求下的更优选择：

\begin{itemize}
    \item 以定点跟踪品质为主、工作点变化小：LQR的最优增益平衡仍有优势；
    \item 需要强未建模动态对消、且能提供高质量角加速度信号：INDI更合适；
    \item 需要工作点鲁棒的交叉项补偿、且希望规避角加速度估计噪声：
          \textbf{级联架构是首选}。
\end{itemize}

\subsection{结论}
\label{sec:cascade_conclusion}

本章提出的级联姿态-角速度控制架构，通过时间尺度分离将姿态调节（外环，
误差四元数P/PI）与交叉项补偿（内环，在线$\mathbf{B}_{\mathrm{true}}$分配）解耦，
在本构型上实现了约两个数量级的耦合抑制改善（巡航工况$0.02$--$0.04^\circ$），
并在tail-sitter万向锁附近保持有效。其优势源于结构解耦而非增益提高：
内环的高带宽使交叉耦合扰动在传播到姿态环之前即被消除。
该架构规避了INDI的角加速度估计噪声、LQR的工作点漂移、SAS的对角近似三大缺陷，
计算开销与INDI相当，是概念阶段向工程实现过渡的推荐控制结构。
严格稳定性证明、参数摄动鲁棒性与执行器饱和下的补偿退化，
是后续工作需要补强的方向。
